% GENERATED FILE. Edit canonical lessons, then run npm run generate:books.
% canonical-source-hash: fnv1a64:41903a3e65d47a56
% canonical-lessons: SA-C260-mam, SA-C260-tvam, SA-C260-singular-object-practice

\chapter{Me and you as objects}
\label{ch:sa-me-and-you-objects}
\hlchaptermodality{260}

\begin{chapteropening}
\textbf{By the end of this chapter:} \emph{I can say that someone sees me or one familiar person addressed.}

Complete chapter 260 by choosing and producing the first- or second-person singular object in a known seeing frame.
\end{chapteropening}

\section[\texorpdfstring{\sk{माम्} (mām)}{माम् (mām)}]{\sk{माम्} (mām) — me, as the object}

\label{lesson:SA-C260-mam}



\begin{quote}
Say \textbf{\sk{सः} \sk{पश्यति}} (\emph{saḥ paśyati}), “he sees”. Keep the seeing verb. Now make the person seen be the speaker.
\end{quote}

\subsection*{You'll want to know: \sk{माम्}}
\textbf{\sk{माम्}} (\emph{mām}) means “me” when I am the object of the action.

\begin{quote}
\textbf{\sk{सः} \sk{माम्} \sk{पश्यति।}}
\end{quote}

\begin{quote}
\emph{saḥ mām paśyati.}
\end{quote}

\begin{quote}
\textbf{He sees me.}
\end{quote}

The verb still says who does the seeing. \textbf{\sk{माम्}} names the person the seeing reaches.

Whitney's \emph{Sanskrit Grammar} §491 lists \textbf{\sk{माम्}} as the independent accusative singular of the first-person pronoun. “Accusative” is the case label; “the object” is the job you can hear in this sentence.

\begin{cousinweb}
As a grammatical form, \textbf{\sk{माम्}} comes from the first-person pronoun paradigm; it is not a word from a separate lexical root. Whitney's \emph{Sanskrit Grammar} §491 records it as the independent accusative singular. This lesson makes no deeper historical derivation claim.

Source: Whitney, \emph{Sanskrit Grammar} (1889), §491.
\end{cousinweb}

\subsection*{Your turn}
Say these aloud:

\begin{itemize}
\raggedright
  \item “me” as the object — \emph{mām}
  \item “he sees me” — \emph{saḥ mām paśyati}
  \item the person seen is the speaker — \emph{mām}
\end{itemize}

\begin{tcolorbox}[breakable,colback=teal!4,colframe=teal!35!black,title={Before you move on}]
Someone sees the speaker. Which one-word object do you use? (\textbf{\sk{माम्}}, \emph{mām}.)
\end{tcolorbox}
\section[\texorpdfstring{\sk{त्वाम्} (tvām)}{त्वाम् (tvām)}]{\sk{त्वाम्} (tvām) — you, one familiar person, as the object}

\label{lesson:SA-C260-tvam}



\begin{quote}
Say “he sees me”: \textbf{\sk{सः} \sk{माम्} \sk{पश्यति}} (\emph{saḥ mām paśyati}). Keep the frame and turn the object toward one familiar person.
\end{quote}

\subsection*{You'll want to know: \sk{त्वाम्}}
\textbf{\sk{त्वाम्}} (\emph{tvām}) means “you” when one familiar addressee is the object.

\begin{quote}
\textbf{\sk{सः} \sk{त्वाम्} \sk{पश्यति।}}
\end{quote}

\begin{quote}
\emph{saḥ tvām paśyati.}
\end{quote}

\begin{quote}
\textbf{He sees you.}
\end{quote}

The opening \textbf{tv-} also appears in familiar \textbf{tvam} and \textbf{tava}. Here the long \textbf{ā} and final \textbf{m} belong to the independent accusative singular. Whitney §491 places \textbf{\sk{त्वाम्}} beside \textbf{\sk{माम्}} in the first- and second-person paradigms.

\begin{cousinweb}
As a grammatical form, \textbf{\sk{त्वाम्}} comes from the second-person pronoun paradigm; it is not a word from a separate lexical root. Whitney's \emph{Sanskrit Grammar} §491 records it as the independent accusative singular. This lesson makes no deeper historical derivation claim.

Source: Whitney, \emph{Sanskrit Grammar} (1889), §491.
\end{cousinweb}

\subsection*{Your turn}
Say these aloud:

\begin{itemize}
\raggedright
  \item “he sees me” — \emph{saḥ mām paśyati}
  \item “he sees you” — \emph{saḥ tvām paśyati}
  \item speaker as object — \emph{mām}; one addressee as object — \emph{tvām}
\end{itemize}

\begin{tcolorbox}[breakable,colback=teal!4,colframe=teal!35!black,title={Before you move on}]
One familiar person is seen. Which object form addresses that person? (\textbf{\sk{त्वाम्}}, \emph{tvām}.)
\end{tcolorbox}
\section[Practice]{Is the one person me or you?}

\label{lesson:SA-C260-singular-object-practice}



\begin{quote}
Give the two singular objects: “me” and “you”. (\emph{Mām; tvām.})
\end{quote}

\begin{grammarlens}[title={Grammar Lens: choose by person}]
There is one person in either sentence. Choose by who that person is:

\begin{itemize}
\raggedright
  \item the speaker as object: \textbf{\sk{माम्}} (\emph{mām});
  \item one familiar addressee as object: \textbf{\sk{त्वाम्}} (\emph{tvām}).
\end{itemize}

Both are independent full forms and both take the same place before \textbf{\sk{पश्यति}}. The choice changes the person seen, not the verb.
\end{grammarlens}

\subsection*{Your turn}
\begin{itemize}
\raggedright
  \item \emph{Say it:} “he sees me” — \emph{saḥ mām paśyati}
  \item \emph{Say it:} “he sees you” — \emph{saḥ tvām paśyati}
  \item \emph{Say it:} the rule — choose \emph{mām} for the speaker, \emph{tvām} for one familiar addressee
  \item \emph{Contrast:} \emph{mām; tvām; mām; tvām}
\end{itemize}

\begin{tcolorbox}[breakable,colback=teal!4,colframe=teal!35!black,title={Before you move on}]
Say both sentences without looking back, then name what changes. (\emph{Saḥ mām paśyati; saḥ tvām paśyati; the person used as the object changes.})
\end{tcolorbox}
